\setcounter{page}{176}
\section{The fundamental local isomorphism}

We have defined duality functors \(f^b\) for finite morphisms and \(f^!\) for projective morphisms in separate settings. In this section we establish a key local isomorphism that unifies these two constructions, laying the groundwork for defining \(f^!\) for arbitrary morphisms of preschemes.

Let \(X = \operatorname{Spec}A\) be an affine scheme. A sequence \(f_1,\dots,f_n\in A\) is called an \textit{$A$-sequence} (regular sequence) if \(f_1\) is a non-zero divisor in \(A\), and for each \(2\le i\le n\), \(f_i\) is a non-zero divisor in \(A/(f_1,\dots,f_{i-1})\).
Let \(J\subseteq A\) be the ideal generated by \(f_1,\dots,f_n\), and let \(\mathcal{J}\) be the corresponding ideal sheaf on \(X\).

We define the \textit{Koszul complex} \(K_\bullet(\underline{f})\) (cf.\ EGA III.I.1, with adjusted sign conventions):
\[
K_{p}(\underline{f})=\Lambda^{p}\bigl(\sheaf{O}_X^{\oplus n}\bigr),\quad p=0,\dots,n .
\]
Let \(e_1,\dots,e_n\) be the standard basis of \(\sheaf{O}_X^{\oplus n}\), and \(\underline{f}=(f_1,\dots,f_n)\in A^n\).

\setcounter{page}{177}
The boundary differential
\[
d_{p}\colon K_{p}(\underline{f}) \to K_{p-1}(\underline{f})
\]
is defined by
\[
d_{p}\bigl(e_{i_1}\wedge\cdots\wedge e_{i_p}\bigr)
=\sum_{j} (-1)^j f_{i_j}\, e_{i_1}\wedge\cdots\wedge\hat{e}_{i_j}\wedge\cdots\wedge e_{i_p}.
\]

For any sheaf \(\sheaf{F}\in\operatorname{Mod}(X)\), define the dual Koszul complex:
\[
K^\bullet(\underline{f};\sheaf{F})
= \RHom\bigl(K_\bullet(\underline{f}),\sheaf{F}\bigr).
\]
A section
\[
\alpha \in \Gamma\big(K^{p}(\underline{f};\sheaf{F})\big)
\]
is determined by its values
\[
\alpha_{i_1\cdots i_p}
=\alpha\bigl(e_{i_1}\wedge\cdots\wedge e_{i_p}\bigr)
\in \Gamma(X,\sheaf{F}).
\]
The coboundary operator is
\[
(d\alpha)_{i_1\cdots i_{p+1}}
=\sum_j (-1)^j f_{i_j}\,\alpha_{i_1\cdots\hat{i}_j\cdots i_{p+1}}.
\]

We denote the cohomology of \(K^\bullet(\underline{f};\sheaf{F})\) by \(H^\bullet(\underline{f};\sheaf{F})\).

Recall [EGA III.1.1] that for an \(A\)-sequence \(\underline{f}\), the Koszul complex \(K_\bullet(\underline{f})\) is a finite locally free resolution of \(\sheaf{O}_X/\mathcal{J}\), via the augmentation \(K_0(\underline{f})=\sheaf{O}_X\to\sheaf{O}_X/\mathcal{J}\).
\setcounter{page}{178}
Hence we have canonical isomorphisms
\[
\RHom\bigl(\sheaf{O}_X/\mathcal{J},\sheaf{F}\bigr)
\cong \RHom\bigl(K_\bullet(\underline{f}),\sheaf{F}\bigr),
\]
which induce
\[
\psi^{i}\colon \mathcal{Ext}_{X}^{i}\bigl(\sheaf{O}_X/\mathcal{J},\sheaf{F}\bigr)
\stackrel{\sim}{\to} H^{i}(\underline{f};\sheaf{F})
\quad (0\le i\le n).
\]

Define the canonical map
\[
\varphi_{\underline{f}}\colon
\mathcal{Ext}_{X}^{n}\bigl(\sheaf{O}_X/\mathcal{J},\sheaf{F}\bigr)
\to \sheaf{F}/\mathcal{J}\sheaf{F}
\]
by composing \(\psi^n\) with the evaluation map sending \(\alpha\in K^n(\underline{f};\sheaf{F})\) to \(\alpha_{1\cdots n}\in\Gamma(X,\sheaf{F})\).
This map is an isomorphism.

More generally, one has for all \(0\le i\le n\):
\[
\mathcal{Ext}_{X}^{i}\bigl(\sheaf{O}_X/\mathcal{J},\sheaf{F}\bigr)
\cong \mathcal{Tor}_{n-i}\bigl(\sheaf{O}_X/\mathcal{J},\sheaf{F}\bigr).
\]
In particular, \(i=0\) recovers
\[
\mathcal{Ext}_{X}^{0}\bigl(\sheaf{O}_X/\mathcal{J},\sheaf{F}\bigr)
=\sheaf{F}\otimes_{\sheaf{O}_X}\sheaf{O}_X/\mathcal{J}
=\sheaf{F}/\mathcal{J}\sheaf{F}.
\]

\setcounter{page}{179}
\begin{lemma}
Let \(X=\operatorname{Spec}A\), let \(\underline{f}=(f_1,\dots,f_n)\), \(\underline{g}=(g_1,\dots,g_n)\) be two \(A\)-sequences generating the same ideal \(J\subseteq A\), and suppose \(g_i=\sum_j c_{ij}f_j\) with \(c_{ij}\in A\).
Then there is a commutative diagram induced by the determinant \(\det(c_{ij})\) on
\[
\mathcal{Ext}_{X}^{n}\bigl(\sheaf{O}_X/J,\sheaf{F}\bigr).
\]
\end{lemma}
\begin{proof}
The change of basis matrix \((c_{ij})\) induces an isomorphism of Koszul complexes \(K_\bullet(\underline{f})\cong K_\bullet(\underline{g})\), whose degree-\(n\) component is multiplication by \(\det(c_{ij})\). The assertion follows immediately.
\end{proof}

\setcounter{page}{180}
\begin{proposition}[Fundamental Local Isomorphism]
Let \(i\colon Y \hookrightarrow X\) be a closed immersion of locally noetherian preschemes, with \(Y\) a locally complete intersection of codimension \(n\) in \(X\). Let \(\sheaf{F}\) be an \(\sheaf{O}_X\)-module.
There exists a natural functorial isomorphism
\[
\varphi\colon
\mathcal{Ext}_{X}^{n}\bigl(\sheaf{O}_X,\sheaf{F}\bigr)
\stackrel{\sim}{\to} \sheaf{F}\otimes_{\sheaf{O}_X}\omega_{Y/X}.
\]
Furthermore, if \(\sheaf{F}\) is \(i^*\)-acyclic, then
\[
\mathcal{Ext}_{X}^{j}\bigl(\sheaf{O}_Y,\sheaf{F}\bigr)=0 \quad \text{for } j\neq n.
\]
\end{proposition}
\begin{proof}
Let \(\mathcal{J}\) be the ideal sheaf of \(Y\) in \(X\). Since \(Y\) is l.c.i., \(\Lambda^n(\mathcal{J}/\mathcal{J}^2)\) is invertible on \(Y\), and
\[
\sheaf{F}\otimes_{\sheaf{O}_X}\omega_{Y/X}
\cong \RHom_Y\bigl(\Lambda^n(\mathcal{J}/\mathcal{J}^2),\sheaf{F}/\mathcal{J}\sheaf{F}\bigr).
\]
Locally we choose an \(A\)-sequence generating \(\mathcal{J}\) and define \(\varphi\) via \(\varphi_{\underline{f}}\). By Lemma 7.1, the isomorphism is independent of local basis choices, hence glues to a global isomorphism.

The condition that \(\sheaf{F}\) is \(i^*\)-acyclic is equivalent to \(\mathcal{Tor}_j(\sheaf{O}_Y,\sheaf{F})=0\) for \(j\neq0\). Using the Ext–Tor duality isomorphism above, this implies vanishing of \(\mathcal{Ext}^j(\sheaf{O}_Y,\sheaf{F})\) for \(j\neq n\).
\end{proof}

\setcounter{page}{181}
\begin{corollary}
Let \(i\colon Y\hookrightarrow X\) be a codimension-\(n\) locally complete intersection closed immersion. There is a natural functorial isomorphism
\[
\eta_i\colon
i^b\bigl(F^\bullet\bigr)
\to \Lfst i^*\bigl(F^\bullet\bigr)\otimes\omega_{Y/X}[-n]
\]
for all \(F^\bullet\in D(X)\).
\end{corollary}
\begin{proof}
We use right/left derived functor conventions; note \(\otimes\) is over \(Y\) since \(\omega_{Y/X}\) is locally free on \(Y\).
If \(F^\bullet\) reduces to a single \(i^*\)-acyclic sheaf \(\sheaf{F}\), the left-hand side is \(\mathcal{Ext}_X^n(\sheaf{O}_Y,\sheaf{F})\) concentrated in degree \(n\), while the right-hand side is isomorphic via Proposition 7.2.

The functor \(i^*\) has finite cohomological dimension, since its derived functors are Tor groups computed by the finite-length Koszul resolution. Thus every complex admits an \(i^*\)-acyclic resolution. Applying [I.7.4] yields the functorial isomorphism.
\end{proof}

\setcounter{page}{182}
\begin{proposition}
\begin{enumerate}
\item[(a)] Let \(Z\xrightarrow{j}Y\xrightarrow{i}X\) be two locally complete intersection closed immersions of codimensions \(m,n\) respectively. There is a commutative diagram of functorial isomorphisms relating \((ij)^b\) with \(j^b,i^b\) via composition of dualizing sheaves and shift functors.

\item[(b)] Let \(i\colon Y\hookrightarrow X\) be an l.c.i.\ closed immersion, \(f\colon X'\to X\) flat, and form the fibre product \(Y'=Y\times_X X'\) with projections \(g\colon Y'\to Y\). Then there is a canonical commutative diagram
\[
\begin{CD}
f^*i^b(F^\bullet) @>{\alpha}>> g^b\Lfst f^*(F^\bullet)
\end{CD}
\]
compatible with base-change isomorphisms for dualizing sheaves.
\end{enumerate}
\end{proposition}
\setcounter{page}{183}
